% Zone „Das kennst du schon“ – aus bank/funktionsscharen-und-ortskurven/zone.jsonl (zusammenbau v0.1)
\einheitenkopf[zone][Kennst du schon]{Das kennst du schon}
\setcounter{aufgabe}{0}

%% TODO Titel: Ich-kann-Satz fehlt in der Bank (Platzhalter: Kettenname) – Nr. 1
%% TODO Anweisung: ein Satz über der ersten Teilaufgabe fehlt in der Bank – Nr. 1
\begin{aufgabe}{Terme mit einem Buchstaben mehr umformen}
\begin{teile}
\teil Klammere $b$ aus: $5b + b^2$. $b\,($ \_\_ $)$
\teil Vereinfache für $t \neq 0$: $\frac{3}{t} \cdot \frac{t^2}{6}$. \leerfeld
\end{teile}
\end{aufgabe}

%% TODO Titel: Ich-kann-Satz fehlt in der Bank (Platzhalter: Kettenname) – Nr. 2
%% TODO Anweisung: ein Satz über der ersten Teilaufgabe fehlt in der Bank – Nr. 2
\begin{aufgabe}{Lineare Gleichungen nach einem beliebigen Buchstaben auflösen}
\begin{teile}
\teil Löse nach $c$ auf: $5 + 2c = 11$. $c =$ \_\_
\teil Löse nach $m$ auf: $\frac{m}{4} - 1 = \frac{1}{2}$. $m =$ \_\_
\end{teile}
\end{aufgabe}

%% TODO Titel: Ich-kann-Satz fehlt in der Bank (Platzhalter: Kettenname) – Nr. 3
%% TODO Anweisung: ein Satz über der ersten Teilaufgabe fehlt in der Bank – Nr. 3
\begin{aufgabe}{Quadratische Gleichungen und Diskriminante}
\begin{teile}
\teil Löse: $x^2 = 25$. $x_1 =$ \_\_, $x_2 =$ \_\_
\teil Wie viele Lösungen hat $x^2 - 4x + 5 = 0$? Nutze die Diskriminante $D = b^2 - 4ac$.
\end{teile}
\end{aufgabe}

%% TODO Titel: Ich-kann-Satz fehlt in der Bank (Platzhalter: Kettenname) – Nr. 4
%% TODO Anweisung: ein Satz über der ersten Teilaufgabe fehlt in der Bank – Nr. 4
\begin{aufgabe}{Parabelparameter kennen}
\begin{teile}
\teil Die Normalparabel wird um 3 nach oben verschoben. Gib die Gleichung an. $y =$ \_\_
\teil Wie unterscheidet sich $y = -\frac{1}{2}x^2$ von der Normalparabel?
\end{teile}
\end{aufgabe}

%% TODO Titel: Ich-kann-Satz fehlt in der Bank (Platzhalter: Kettenname) – Nr. 5
%% TODO Anweisung: ein Satz über der ersten Teilaufgabe fehlt in der Bank – Nr. 5
\begin{aufgabe}{Ableitungsregeln sicher anwenden, auch mit Vorfaktoren}
\begin{teile}
\teil Leite ab: $f(x) = 5x^3$. $f'(x) =$ \_\_
\teil Leite ab: $f(x) = \frac{1}{6}x^3 + 2x$. $f'(x) =$ \_\_
\end{teile}
\end{aufgabe}

%% TODO Titel: Ich-kann-Satz fehlt in der Bank (Platzhalter: Kettenname) – Nr. 6
%% TODO Anweisung: ein Satz über der ersten Teilaufgabe fehlt in der Bank – Nr. 6
\begin{aufgabe}{Extrem- und Wendepunktkriterien an fester Funktion führen}
\begin{teile}
\teil Bestimme die Extremstelle von $f(x) = x^2 - 8x$ und ihre Art. $x =$ \_\_
\teil Bestimme die Wendestelle von $f(x) = x^3 - 6x^2 + 1$. $x =$ \_\_
\end{teile}
\end{aufgabe}

%% TODO Titel: Ich-kann-Satz fehlt in der Bank (Platzhalter: Kettenname) – Nr. 7
%% TODO Anweisung: ein Satz über der ersten Teilaufgabe fehlt in der Bank – Nr. 7
\begin{aufgabe}{Stammfunktion bilden und bestimmtes Integral auswerten}
\begin{teile}
\teil Gib eine Stammfunktion an: $f(x) = 6x^2$. $F(x) =$ \_\_
\teil Berechne das Integral von $1$ bis $3$ über $\frac{1}{2}x$. \leerfeld
\end{teile}
\end{aufgabe}

%% TODO Titel: Ich-kann-Satz fehlt in der Bank (Platzhalter: Kettenname) – Nr. 8
%% TODO Anweisung: ein Satz über der ersten Teilaufgabe fehlt in der Bank – Nr. 8
\begin{aufgabe}{Terme mit einem Buchstaben mehr umformen – Fehler finden}
\begin{teile}
\teil Finde den Fehler und rechne richtig: $e^{4m} \cdot e^{-m} = e^{-4m^2}$.
\end{teile}
\end{aufgabe}

%% TODO Titel: Ich-kann-Satz fehlt in der Bank (Platzhalter: Kettenname) – Nr. 9
%% TODO Anweisung: ein Satz über der ersten Teilaufgabe fehlt in der Bank – Nr. 9
\begin{aufgabe}{Terme mit einem Buchstaben mehr umformen – selbst rechnen}
\begin{teile}
\teil Vereinfache: $\frac{e^{3r}}{e^{-r}}$. \leerfeld
\end{teile}
\end{aufgabe}
